Niniejsza sekcja opisuje metodologię zastosowaną w opracowaniu i ocenie modeli rozpoznawania twarzy z naciskiem na niezawodność wobec okluzji. Szczegółowo omawiamy przygotowanie danych, architektury modeli, procedury treningowe i protokoły ewaluacji.

\subsection{Przygotowanie danych i przetwarzanie wstępne}
\subsubsection{Przegląd datasetu}

\paragraph{}Badania przeprowadzono na zbiorze danych CASIA-WebFace, zawierającym ponad 494 tys. kolorowych fotografii twarzy, należących do 10,575 unikatowych tożsamości. Dane podzielono na dwa rozłączne podzbiory: treningowy (90\%) i testowy (10\%). Dla monitorowania zbieżności modelu wykorzystywano stały zestaw 500 par (walidacja) z części testowej, zgodnie z pipeline’em treningowym.

\subsubsection{Wyrównanie na podstawie punktów orientacyjnych}

\paragraph{}Wszystkie obrazy są wstępnie przetwarzane za pomocą modelu RetinaFace od InsightFace w celu wykrycia 5 punktów orientacyjnych twarzy: lewe oko, prawe oko, czubek nosa i kąciki ust. Te punkty orientacyjne umożliwiają konsystentną normalizację twarzy przy użyciu transformacji podobieństwa do standardowej rozdzielczości 112×112. Większość obrazów ma towarzyszący plik JSON zawierający:
\begin{itemize}
    \item Współrzędne ramki ograniczającej: \texttt{bbox} = $[x_1, y_1, x_2, y_2]$
    \item 5 punktów orientacyjnych: \texttt{landmarks} = $\{$lewe\_oko, prawe\_oko, nos, lewy\_ust, prawy\_ust$\}$
\end{itemize}

W przypadku braku JSON, zastosowano rezerwowe skalowanie obrazu do 112×112. Brak pliku JSON wynikał z niepowodzenia detekcji punktów orientacyjnych dla niskiej jakości lub nietypowych obrazów (ok. 2,5\% zdjęć).

Wyrównanie twarzy jest wykonywane przy użyciu referencyjnej konfiguracji punktów orientacyjnych ArcFace:
\begin{equation}
M = \text{Similarity Transform}(\text{wykryte\_punkty}, \text{punkty\_referencyjne})
\end{equation}

gdzie punkty referencyjne dla obrazów 112×112 są ustalone na:
\begin{equation}
\text{arcface\_dest} = \begin{bmatrix} 38.29 & 51.70 \\ 73.53 & 51.50 \\ 56.03 & 71.74 \\ 41.55 & 92.37 \\ 70.73 & 92.20 \end{bmatrix}
\end{equation}

Punkty te stanowią standardową konfigurację stosowaną w ArcFace i zostały zaczerpnięte z dokumentacji kodu oryginalnego modelu, gdzie określono je jako referencyjne pozycje cech twarzy w wyrównanych obrazach.

Ta procedura zapewniła konsystentne wyrównanie zgodne z oryginalnym modelem Buffalo\_L (ArcFace Large), co jest kluczowe dla transferu wag.

\subsubsection{Strategia augmentacji danych}
\label{subsec:augmentation}

Wszystkie modele są trenowane z syntetyczną okluzją w celu symulacji zasłonięć w regionie oczu. Implementacja stosuje poziomy pasek okluzji o stałej wysokości i losowym kolorze.

\begin{itemize}
    \item \textbf{Wysokość}: 20 pikseli
    \item \textbf{Szerokość}: pełna szerokość obrazu (0 do 112)
    \item \textbf{Pozycja pionowa}: $y = 52 \pm 5$ pikseli (region oczu)
    \item \textbf{Kolor}: losowe RGB $(r, g, b)$, $r,g,b \in [0, 255]$
\end{itemize}

Prawdopodobieństwo wystąpienia okluzji różni się między wariantami zgodnie z kodem treningowym:
\begin{itemize}
    \item \textbf{Aligned\_Pretrained}: $p=0.7$
    \item \textbf{Aligned\_Pretrained\_Transformers}: $p=0.7$
    \item \textbf{CBAM\_Block}: $p=0.5$
\end{itemize}

W fazie testowej kolor okluzji jest stały (czarny), co zapewnia powtarzalność i porównywalność wyników pomiędzy modelami.

\subsection{Ewaluacja modeli typu SOTA}
W celu zapewnienia rzetelnego punktu odniesienia oraz właściwej kontekstualizacji uzyskanych wyników przeprowadzono ewaluację wybranych modeli typu state-of-the-art (SOTA), w tym modeli dostępnych w bibliotece InsightFace oraz modelu referencyjnego.

Wybrane modele obejmowały zarówno rozwiązania o wysokiej dokładności, jak i lekkie architektury zoptymalizowane pod kątem wydajności obliczeniowej, a także model referencyjny starszej generacji.

Testowane modele:
\begin{itemize}
\item ArcFace Large (buffalo\_l) – model oparty na architekturze IResNet-50, dostępny w bibliotece InsightFace,
\item ArcFace Small (buffalo\_s) – lekka architektura oparta na MobileFaceNet, przeznaczona do zastosowań mobilnych, dostępna w bibliotece InsightFace,
\item VGGFace2 (ResNet-50) – model referencyjny starszej generacji (2018 r.), zaimplementowany z wykorzystaniem biblioteki \textit{keras-vggface}, niebędący częścią frameworka InsightFace.
\end{itemize}

Wszystkie modele poddano temu samemu protokołowi identyfikacji (galeria–sonda) oraz symulacji okluzji w postaci poziomego zasłonięcia regionu oczu o wysokości 20 pikseli. Należy jednak podkreślić, że część modeli SOTA wykorzystuje własny pipeline ekstrakcji cech: dla ArcFace (InsightFace) embeddingi są wyznaczane na podstawie pełnego obrazu z detekcją twarzy, natomiast dla VGGFace (ResNet50) stosowane jest przycięcie względem \texttt{bbox} i skalowanie do 224×224. Zostało to utrzymane zgodnie z implementacjami referencyjnymi w repozytorium. Szczegółowy opis protokołu ewaluacji przedstawiono w dalszej części dokumentu.

\subsection{Model bazowy}
\subsubsection{Wybór IResNet50}
Wybraliśmy IResNet50 (Ulepszona sieć ResNet50) jako architekturę sieci głównej ze względu na jej sprawdzoną wydajność w modelu ArcFace-Large (buffalo\_l) z projektu InsightFace. Implementacja w repozytorium zawiera warstwy BN i PReLU, a warstwa \texttt{features} (BN1d) ma zamrożone wagi (\texttt{requires\_grad=False}), co ogranicza adaptację końcowego wektora cech w trakcie fine-tuningu. IResNet50 składa się z:
\begin{itemize}
    \item Wstępna konwolucja: $3 \times 3$ conv, 64 kanały, stride 1
    \item Warstwy 1-4: Bloki resztkowe o różnych głębokościach
    \item Normalizacja wsadowa i aktywacje PReLU na całej sieci
    \item Wymiar cechy: osadzania 512-wymiarowe
\end{itemize}

\subsubsection{Hybrydowa strategia transferu wag}
Bezpośredni transfer z modelu buffalo\_l opartego na ONNX do PyTorch okazał się trudny ze względu na niezgodności architektoniczne. Zastosowano hybrydowe dopasowanie wag (dopasowanie nazw + dopasowanie kształtu), z reinicjalizacją warstw bez dopasowania. Proces ten jest realizowany przez \texttt{load.py} i zapewnia stabilny punkt startowy dla fine-tuningu.

\paragraph{Pomyślne transfery:} $\approx 90\%$ wag bloków ResNet zostało pomyślnie zmapowanych, w tym:
\begin{itemize}
    \item Warstwy Conv2d w ścieżkach resztkowych
    \item Większość parametrów BatchNorm (wagi, obciążenia, statystyki uruchamiające)
    \item Wagi PReLU (gdzie konwencje nazewnictwa się zgadzały)
\end{itemize}

\paragraph{Warstwy reinicjalizowane:} 26 warstw wymagało inicjalizacji Kaiming Normal (He) ze względu na niezgodność architektoniczną:
\begin{itemize}
    \item layer1.2.bn3: statystyki BatchNorm3D (running\_mean, running\_var)
    \item layer2.3.bn2, layer2.3.bn3: parametry BatchNorm
    \item layer2.3.prelu: waga PReLU (parametr $\alpha$)
    \item layer3.12-13.bn2-3: podobne instancje BatchNorm i PReLU
\end{itemize}

Formuła inicjalizacji:
\begin{equation}
w \sim \mathcal{N}(0, \sqrt{\frac{2}{n_{in}(1 + a^2)}})
\end{equation}
gdzie $a$ jest ujemnym nachyleniem (0 dla ReLU, 0.25 dla PReLU w naszym przypadku - później trenowane).

\subsection{Funkcje straty}
\label{subsec:loss}

\subsubsection{Strata ArcMargin}
Podstawową funkcją straty w modelu bazowym oraz w wariancie CBAM\_Block jest strata ArcMargin, marginalna strata softmax, która wymusza separowalność kątową w przestrzeni osadzania. W implementacji przyjęto parametry $m=0.5$ i $s=30.0$ (zbieżne z ArcFace).

W treningu optymalizowana jest standardowa \textit{CrossEntropy} na wyjściu ArcMargin:
\begin{equation}
\mathcal{L}_{\text{ArcFace}} = \text{CE}(\text{ArcMargin}(\mathbf{e}), y)
\end{equation}

\subsubsection{Strata złożona dla wariantu Transformer}
W wariancie Transformer zastosowano funkcję straty łączoną, gdzie gałąź podstawowa (ArcMargin) jest wzbogacana o pomocniczą stratę klasyfikacji z transformera (CrossEntropy):
\begin{equation}
\mathcal{L}_{\text{total}} = (1-\alpha)\,\mathcal{L}_{\text{ArcFace}} + \alpha\,\mathcal{L}_{\text{Trans}}
\end{equation}
gdzie $\alpha = 0.4$ zgodnie z parametrami w kodzie treningu.

\subsection{Modyfikacja Architektury: Strategie Odporności na Okluzję}
Zaproponowane modyfikacje architektury ograniczono do wariantów faktycznie obecnych w implementacji. Wszystkie modele bazują na IResNet50 i różnią się wyłącznie dodatkowymi mechanizmami odporności na okluzję.

\paragraph{Wariant A: Model bazowy}
Wariant odniesienia wykorzystuje standardowy backbone IResNet50 oraz głowicę ArcMargin do klasyfikacji tożsamości. Nie stosuje mechanizmów uwagi ani dodatkowych głowic pomocniczych. Jego rolą jest stanowić stabilny punkt odniesienia do oceny wpływu modyfikacji architektury.

\paragraph{Wariant B: CBAM zintegrowany w blokach rezydualnych}
W tym wariancie moduły CBAM umieszczono wewnątrz każdego bloku rezydualnego (CBAMBasicBlock), bezpośrednio przed połączeniem rezydualnym. Zastosowano uwagę kanałową i przestrzenną, co pozwala na selektywne wzmacnianie istotnych map cech na wielu poziomach abstrakcji. Celem wariantu jest uzyskanie odporności na okluzję poprzez rozproszoną filtrację zakłóceń.

\paragraph{Wariant C: Gałąź pomocnicza Transformer}
Wariant ten modyfikuje backbone tak, aby zwracał mapę cech o rozmiarze $7\times7$, a następnie dodaje pomocniczą gałąź Transformer (Branch-2) uczoną równolegle z głowicą ArcMargin (Branch-1). W trakcie treningu wykorzystywana jest strata łączona (ArcMargin + CrossEntropy z gałęzi Transformer), natomiast podczas ewaluacji wykorzystywana jest wyłącznie gałąź standardowa.

\subsubsection{Strategia uczenia i stabilizacji}
We wszystkich wariantach zastosowano transfer learning z wagami wstępnie wytrenowanymi, a w wariancie CBAM\_Block dodatkowo etapowe zamrażanie/odmrażanie wag w celu stabilizacji uczenia. Augmentacja danych obejmowała syntetyczną okluzję regionu oczu, spójną z protokołem ewaluacyjnym.


\subsection{Procedura Treningu}
\subsubsection{Model bazowy (Aligned\_Pretrained)}
Model bazowy trenowany jest jako klasyfikator z głowicą ArcMargin.

\begin{table}[h]
\centering
\small
\begin{tabular}{ll}
\toprule
\textbf{Hiperparametr} & \textbf{Wartość} \\
\midrule
Rozmiar partii & 32 \\
Epoki & 25 \\
Optymalizator & SGD (momentum=0.9) \\
LR & $1 \times 10^{-2}$ \\
Zanik wag & brak (nie ustawiono) \\
Wczesne zatrzymanie & Tak (cierpliwość=5) \\
Walidacja & 500 par z testu (średnia cos) \\
Augmentacja okluzji & $p=0.7$, wysokość 20 px \\
\bottomrule
\end{tabular}
\caption{Hiperparametry modelu bazowego (Aligned\_Pretrained).}
\label{tab:baseline_config}
\end{table}

\subsubsection{Gałąź pomocnicza Transformer}
Model z transformatorem wykorzystuje złożoną funkcję straty, łącząc ArcMargin i klasyfikację gałęzi pomocniczej.

\begin{table}[h]
\centering
\small
\begin{tabular}{ll}
\toprule
\textbf{Hiperparametr} & \textbf{Wartość} \\
\midrule
Rozmiar partii & 32 \\
Epoki & 25 \\
Optymalizator & SGD (momentum=0.9, weight\_decay=$5\times 10^{-4}$) \\
LR & $1 \times 10^{-2}$ \\
ALPHA (balans strat) & 0.4 \\
Transformer layers/heads & 6 / 8 \\
Wczesne zatrzymanie & Tak (cierpliwość=5) \\
Walidacja & 500 par z testu (średnia cos) \\
Augmentacja okluzji & $p=0.7$, wysokość 20 px \\
\bottomrule
\end{tabular}
\caption{Hiperparametry modelu z gałęzią Transformer.}
\label{tab:transformer_config}
\end{table}

\subsubsection{CBAM\_Block (uwaga w każdym bloku)}
Trening realizowany jest dwuetapowo z zamrożeniem backbone'u w pierwszej epoce.

\begin{table}[h]
\centering
\small
\begin{tabular}{ll}
\toprule
\textbf{Hiperparametr} & \textbf{Wartość} \\
\midrule
Rozmiar partii & 64 \\
Epoki & 30 \\
Optymalizator & SGD (momentum=0.9, weight\_decay=$5\times 10^{-4}$) \\
LR start & 0.1 \\
Harmonogram LR & StepLR (step\_size=10, gamma=0.1) \\
Zamrażanie & Epoka 1: tylko CBAM + head \\
Walidacja & 500 par z testu (średnia cos) \\
Augmentacja okluzji & $p=0.5$, wysokość 20 px \\
\bottomrule
\end{tabular}
\caption{Hiperparametry modelu CBAM\_Block.}
\label{tab:cbam_block_config}
\end{table}

\subsection{Protokół ewaluacji}
\subsubsection{Konstrukcja zestawu testowego}
Protokół ewaluacji stosuje zadanie identyfikacji galeria-vs-sonda:

\paragraph{Zestaw galerii:} Skonstruowany z pierwszej połowy sesji (folderów) dla każdej osoby w podziale testowym. Dla każdej osoby osadzania z wielu czystych (nie okluzjowanych) obrazów są uśredniane:

\begin{equation}
\mathbf{e}_{\text{gallery}} = \frac{1}{n} \sum_{i=1}^{n} \text{L2-Normalize}(\mathbf{e}_i)
\end{equation}

Zapewnia to niezawodny prototyp osoby do dopasowania podobieństwa.

\paragraph{Zestaw sondy:} Skonstruowany z drugiej połowy sesji. Te obrazy są poddane syntetycznej okluzji podczas ewaluacji (szczegóły w sekcji \ref{subsubsec:eval_occlusion}).

\subsubsection{Ekstrakcja embeddingu i indeksowanie}
\begin{enumerate}
    \item Ekstrahuj embedding dla wszystkich obrazów galerii i sondy
    \item Normalizuj L2: $\hat{\mathbf{e}} = \frac{\mathbf{e}}{\|\mathbf{e}\|_2}$
    \item Indeksuj embedding galerii przy użyciu FAISS IndexFlatIP (Inner Product), gdzie podobieństwo cosinusowe na wektorach znormalizowanych L2 równa się iloczynowi wewnętrznemu
\end{enumerate}

\paragraph{Rozdział walidacji i ewaluacji:} w trakcie treningu stosowana jest pomocnicza walidacja na 500 parach (średnia cosinusowa czyste vs okluzjowane), próbkowanych z części testowej. Wyniki Rank-1/Rank-3 pochodzą z niezależnego pipeline’u FAISS dla zadania identyfikacji galeria-vs-sonda.

\subsubsection{Symulacja okluzji w ewaluacji}
\label{subsubsec:eval_occlusion}

Podczas ewaluacji każdy obraz sondy jest syntetycznie okluzjowany przy użyciu deterministycznego protokołu, aby zapewnić powtarzalność:

\begin{itemize}
    \item \textbf{Wysokość okluzji}: 20 pikseli
    \item \textbf{Szerokość okluzji}: jeśli dostępny \texttt{bbox} z adnotacji, pasek ograniczany jest do zakresu $[x_1, x_2]$; w przeciwnym razie pełna szerokość obrazu (0 do 112)
    \item \textbf{Kolor}: Czarny (0,0,0) [deterministyczny, w przeciwieństwie do losowej augmentacji treningowej]
    \item \textbf{Pozycja pionowa}: Region oczu, obliczony z punktów orientacyjnych:
    \begin{equation}
    c_y = \frac{y_{\text{lewe\_oko}} + y_{\text{prawe\_oko}}}{2}
    \end{equation}
    \item \textbf{Fallback}: Jeśli punkty orientacyjne niedostępne, $c_y = \frac{h}{2} - 10 = 46$ (dla 112×112)
    \item \textbf{Zakres Y}: $[c_y - 10, c_y + 10]$
\end{itemize}

Deterministyczna czarna okluzja (w przeciwieństwie do losowych kolorów w treningu) zapewnia rygorystyczną, powtarzalną ewaluację podczas testowania uogólnienia na wariacje kolorów.

\paragraph{Wyrównanie w ewaluacji:} obrazy testowe w zbiorze \texttt{webface\_112x112} są już wstępnie wyrównane.

\subsubsection{Wyszukiwanie podobieństwa i metryki}
Dla każdego obrazu sondy $p$:
\begin{enumerate}
    \item Zapytaj indeks FAISS z $k=3$ najbliższymi sąsiadami
    \item Pobierz top-3 tożsamości galerii ze wynikami podobieństwa: $(s_1, s_2, s_3)$
\end{enumerate}

\paragraph{Dokładność Rank-k:} 
\begin{equation}
\text{Rank-}k = \frac{\text{\# zapytań z prawidłową tożsamością w top-}k}{N_{\text{ogółem}} }
\end{equation}

W szczególności:
\begin{itemize}
    \item \textbf{Dokładność Rank-1}: $P(\text{tożsamość}[0] = \text{rzeczywistość})$
    \item \textbf{Dokładność Rank-3}: $P(\text{rzeczywistość} \in \{\text{tożsamość}[0,1,2]\})$
\end{itemize}

\subsubsection{Wyjście i logowanie}
Szczegółowe wyniki są logowane w formacie CSV z kolumnami:
\begin{lstlisting}
query_id, top1_id, sim1, top2_id, sim2, top3_id, sim3, found_in_top3
\end{lstlisting}

\subsection{Środowisko obliczeniowe i implementacja}
\paragraph{}Wszystkie eksperymenty zostały przeprowadzone w środowisku Python z wykorzystaniem PyTorch do trenowania modeli, InsightFace do detekcji punktów orientacyjnych oraz FAISS do szybkiego wyszukiwania podobieństwa w przestrzeni embedding'ów. Dla modeli autorskich zastosowano wspólny pipeline danych (wyrównanie, augmentacja, normalizacja, ekstrakcja embedding'ów), co minimalizuje ryzyko różnic wynikających z niejednorodnych procedur wejściowych. Modele SOTA korzystały z własnych, referencyjnych pipeline’ów zgodnie z implementacjami w repozytorium.

Modelowanie architektur odbywało się w tym samym szkielecie IResNet50, a modyfikacje ograniczano do ściśle kontrolowanych miejsc (gałąź pomocnicza Transformer lub moduły CBAM). Taka polityka zmian pozwala przypisać obserwowane różnice wydajności wyłącznie wprowadzonym mechanizmom odporności na okluzję.

\subsection{Dobór hiperparametrów i stabilność uczenia}
\paragraph{}Wybór hiperparametrów oparto na konserwatywnym dostrajaniu: najpierw ustalono stabilny punkt odniesienia (baseline), następnie wprowadzono modyfikacje architektury (CBAM\_Block i Transformer) przy zachowaniu porównywalnego pipeline’u danych. Dla wariantu CBAM\_Block stosowano harmonogram StepLR oraz etapowe zamrażanie/odmrażanie warstw, a dla pozostałych wariantów wczesne zatrzymanie na podstawie walidacji (500 par) w celu ograniczenia niestabilności numerycznej.

Parametry uczenia (rozmiar partii, epoki, współczynniki uczenia) dobrano empirycznie na podstawie obserwowanej zbieżności funkcji straty i stabilności średniej cosinusowej na zbiorze walidacyjnym (500 par z części testowej). Jednolity protokół augmentacji (okluzja w rejonie oczu) utrzymano we wszystkich wariantach, co umożliwia bezpośrednie porównanie modeli.

\subsection{Reprodukowalność i logowanie}
\paragraph{}W kodzie nie ustawiono stałych ziaren losowości, dlatego przebiegi mogą wykazywać niewielką wariancję między uruchomieniami. Dla każdego wariantu zachowywano najlepszy checkpoint według średniej cosinusowej na walidacji (500 par z części testowej), a końcowa ewaluacja wykonywana była wyłącznie na ustalonym zbiorze testowym z deterministyczną okluzją. Wyniki zbiorcze zostały zapisane w postaci tabel zbiorczych, co ułatwia replikację i weryfikację rezultatów.